% ##############################################################################
\section{Calibration}
\label{sec:calibration}

A generated question can be ambiguous, can depend on an output format that some
providers do not follow, or can be answered correctly by weak and strong models
alike. Using such a question as evidence would make the checker's own error
rate unknown. Calibration therefore measures each candidate against reference
targets before the prober may use it, and it keeps only the candidates that
separate a healthy deployment from a weaker one.

% ==============================================================================
\subsection{Protocol and decision rule}
\label{sec:calibration_rule}

A \emph{target} is a pair of a provider slug and a model identifier with a
role. Calibration uses a set $T_H$ of healthy targets, for example the same
model on several providers, and exactly one weak target $t_w$, for example a
smaller model of the same family. The MVP example in the prober documentation
uses a Llama 3.3 70B Instruct model on two providers as healthy targets and a
Llama 3.1 8B Instruct model on a third provider as the weak target. Every
canary is sent $R$ times to every target at temperature $0$. Each reply
receives a grade by \eqref{eq:grade}.

For a canary $c$, let $O_H(c)$ be the $R\,|T_H|$ observations on healthy
targets and $O_W(c)$ the $R$ observations on the weak target. The canary is
\emph{usable} if all of these observations exist, every provider status is
\texttt{ok}, and no grade is $\bot$. For a usable canary define the healthy and
weak accuracies
\begin{equation}
  \label{eq:calibration_accuracy}
  \acc_H(c) = \frac{1}{|O_H(c)|}\sum_{o \in O_H(c)} \mathbb{1}[o\ \text{correct}],
  \qquad
  \acc_W(c) = \frac{1}{|O_W(c)|}\sum_{o \in O_W(c)} \mathbb{1}[o\ \text{correct}].
\end{equation}
The decision for canary $c$ is
\begin{equation}
  \label{eq:calibration_decision}
  \text{accept}(c) \iff \text{usable}(c) \;\wedge\; \acc_H(c) \geq \tau_h
    \;\wedge\; \acc_W(c) \leq \tau_w,
\end{equation}
with defaults $R = 3$, $\tau_h = 1$, and $\tau_w = 0.6$. A provider failure, a
grader failure, or a missing attempt makes the canary unusable. Such a
canary is reported as \texttt{INCOMPLETE\_OBSERVATIONS} and is never
interpreted as evidence about model capability. A usable canary that fails
the healthy test is reported as \texttt{HEALTHY\_ACCURACY\_TOO\_LOW}, and one
that fails the weak test as \texttt{WEAK\_ACCURACY\_TOO\_HIGH}. Figure~\ref{fig:calibration-decision}
shows the decision order, and Algorithm~\ref{alg:calibration} gives the
procedure. Table~\ref{tab:calibration_params} lists the parameters.

% rendered_images:begin
% ```graphviz
% digraph CalibrationDecision {
%   // input : orange | process : blue | accepted : sage green | rejected : coral | decision : white
%   bgcolor="transparent";
%   pad="0.10";
%   splines=spline;
%   nodesep=0.25;
%   ranksep=0.38;
%   node [shape=box, style="rounded,filled", fillcolor="#FFFFFF", color="#9AA9B8", penwidth=1.8, fontname="Helvetica", fontcolor="#243B53", fontsize=11, margin="0.14,0.08", height=0.42];
%   edge [color="#A3B1C0", penwidth=1.3, arrowhead=vee, arrowsize=0.75, fontname="Helvetica", fontsize=9.5, fontcolor="#7B8794"];
%   rankdir=TB;
%   Cand  [label=<<b>Candidate canary c</b><br/><font point-size="9" color="#6B4517">from the bank</font>>, fillcolor="#FBEBD4", color="#D9A85F", fontcolor="#6B4517"];
%   Run   [label=<<b>Run R attempts on every target</b><br/><font point-size="9" color="#41719C">healthy targets and one weak target, grade each reply</font>>, fillcolor="#D3E3F3", color="#7CA6CE", fontcolor="#1F4E79"];
%   D1 [label="All attempts present,\nprovider ok, gradable?"];
%   D2 [label=<Healthy accuracy<br/>≥ τ<SUB>h</SUB>?>];
%   D3 [label=<Weak accuracy<br/>≤ τ<SUB>w</SUB>?>];
%   Acc [label=<<b>CALIBRATED</b><br/><font point-size="9" color="#4F7A5A">persist both accuracies</font>>, fillcolor="#DFEDE0", color="#8FB79A", fontcolor="#2E5A3D"];
%   R1 [label=<<b>INCOMPLETE_OBSERVATIONS</b>>, fillcolor="#F6C6C6", color="#D98C8C", fontcolor="#6B2A2A"];
%   R2 [label=<<b>HEALTHY_ACCURACY_TOO_LOW</b>>, fillcolor="#F6C6C6", color="#D98C8C", fontcolor="#6B2A2A"];
%   R3 [label=<<b>WEAK_ACCURACY_TOO_HIGH</b>>, fillcolor="#F6C6C6", color="#D98C8C", fontcolor="#6B2A2A"];
%   Draft [label=<<b>Remains DRAFT</b><br/><font point-size="9" color="#6B2A2A">never scheduled</font>>, fillcolor="#F6C6C6", color="#D98C8C", fontcolor="#6B2A2A"];
%   Count [label=<At least m canaries<br/>CALIBRATED? (m = 150)>];
%   Act [label=<<b>Activate bank</b><br/><font point-size="9" color="#4F7A5A">CALIBRATED becomes ACTIVE</font>>, fillcolor="#DFEDE0", color="#8FB79A", fontcolor="#2E5A3D"];
%   NoAct [label=<<b>Bank not activated</b><br/><font point-size="9" color="#6B2A2A">CLI exits with code 2</font>>, fillcolor="#F6C6C6", color="#D98C8C", fontcolor="#6B2A2A"];
% 
%   Cand -> Run;
%   Run -> D1;
%   D1 -> D2 [label="Yes"];
%   D1 -> R1 [label="No"];
%   D2 -> D3 [label="Yes"];
%   D2 -> R2 [label="No"];
%   D3 -> Acc [label="Yes"];
%   D3 -> R3 [label="No"];
%   R1 -> Draft [style=dashed];
%   R2 -> Draft [style=dashed];
%   R3 -> Draft [style=dashed];
%   Acc -> Count [label="All candidates evaluated"];
%   Count -> Act [label="Yes"];
%   Count -> NoAct [label="No"];
%   { rank=same; D1; R1; }
%   { rank=same; D2; R2; }
%   { rank=same; D3; R3; }
% }
% ```
% label=fig:calibration-decision
% caption=Decision flow for one candidate canary and for the bank. Blue is a process, green is an accepted outcome, coral is a rejected outcome, and white boxes are decisions.
% rendered_images:end
% render_images:begin
\begin{figure}[H]
  \includegraphics[width=\linewidth]{05_calibration.tex.figs/05_calibration.1.png}
  \caption{Decision flow for one candidate canary and for the bank. Blue is a process, green is an accepted outcome, coral is a rejected outcome, and white boxes are decisions.}
  \label{fig:calibration-decision}
\end{figure}
% render_images:end

\begin{algorithm}[htbp]
  \caption{Calibration of a candidate bank (\texttt{execute\_calibration} and
  \texttt{evaluate\_canaries}).}
  \label{alg:calibration}
  \begin{algorithmic}[1]
    \Require bank $\bank$; healthy targets $T_H$; weak target $t_w$; repetitions
      $R$; thresholds $\tau_h, \tau_w$; concurrency limit $K$
    \Ensure one decision per canary, and a summary
    \State validate: one bank version, unique canary identifiers, unique
      targets, $|T_H| \geq 1$, exactly one weak target
    \ForAll{$(c, t, r) \in \bank \times (T_H \cup \{t_w\}) \times \{1,\dots,R\}$}
      \Comment{at most $K$ calls in flight}
      \State $\mathit{res} \gets$ call target $t$ with the prompt of $c$ at
        temperature $0$
      \If{$\mathit{res}$ has status \texttt{ok}}
        \State $g \gets \grade(c, \mathit{res}.\mathit{text})$
      \Else
        \State $g \gets \bot$
      \EndIf
      \State store the observation $(c, t, r, \mathit{res}.\mathit{status}, g,
        \mathit{latency}, \mathit{cost}, \mathit{text})$
    \EndFor
    \ForAll{$c \in \bank$ in identifier order}
      \If{$c$ is not usable}
        \State reason $\gets$ \texttt{INCOMPLETE\_OBSERVATIONS}
      \ElsIf{$\acc_H(c) < \tau_h$}
        \State reason $\gets$ \texttt{HEALTHY\_ACCURACY\_TOO\_LOW}
      \ElsIf{$\acc_W(c) > \tau_w$}
        \State reason $\gets$ \texttt{WEAK\_ACCURACY\_TOO\_HIGH}
      \Else
        \State reason $\gets$ \texttt{CALIBRATED}
      \EndIf
    \EndFor
    \State \Return the decisions and the per-target summary
  \end{algorithmic}
\end{algorithm}

\begin{table}[htbp]
  \centering
  \small
  \begin{tabularx}{\linewidth}{@{}p{0.34\linewidth}p{0.12\linewidth}X@{}}
    \toprule
    Parameter & Default & Meaning \\
    \midrule
    Repetitions $R$ & 3 & Attempts per canary and target \\
    Healthy threshold $\tau_h$ & 1.0 & Required accuracy over all healthy
      attempts \\
    Weak threshold $\tau_w$ & 0.6 & Largest accepted weak accuracy \\
    Maximum tokens & 64 & Output limit of every call \\
    Timeout & 15 s & Per-call timeout \\
    Concurrency $K$ & 5 & Calls in flight \\
    Minimum survivors $m$ & 150 & Accepted canaries needed to activate the bank
      \\
    \bottomrule
  \end{tabularx}
  \caption{Calibration policy (\texttt{CalibrationPolicy} and the
  \texttt{calibrate} command).}
  \label{tab:calibration_params}
\end{table}

The number of paid provider calls is $|\bank|\,(|T_H| + 1)\,R$. For the example
targets above, this is $300 \times 3 \times 3 = 2{,}700$ calls. The
\texttt{calibrate} command prints this number and stops unless
\texttt{--execute} is passed.

% ==============================================================================
\subsection{Analysis of the acceptance rule}
\label{sec:calibration_analysis}

Suppose that, for a fixed canary, every attempt on a healthy target is correct
with probability $p$ and every attempt on the weak target is correct with
probability $w$, and that attempts are independent. Under these assumptions
the survival probabilities have closed forms.

\begin{proposition}
\label{prop:acceptance}
With $\tau_h = 1$, $R = 3$, and $\tau_w = 0.6$, a usable canary is accepted
with probability $P_H(p)\,P_W(w)$, where
\begin{equation}
  \label{eq:acceptance}
  P_H(p) = p^{\,3|T_H|}, \qquad P_W(w) = (1-w)^3 + 3w(1-w)^2.
\end{equation}
\end{proposition}

\begin{proof}
The condition $\acc_H(c) \geq 1$ requires all $3|T_H|$ healthy attempts to be
correct, which has probability $p^{3|T_H|}$. The weak accuracy takes values in
$\{0, 1/3, 2/3, 1\}$, and $\acc_W(c) \leq 0.6$ holds if and only if at most one of
the three attempts is correct, because $2/3 > 0.6$. The probability of zero or
one successes in three Bernoulli trials is $(1-w)^3 + 3w(1-w)^2$.
\end{proof}

Figure~\ref{fig:calibration-acceptance} plots both factors, and
Table~\ref{tab:acceptance_values} lists values. Three observations follow.

\begin{figure}[htbp]
  \centering
  \includegraphics[width=\linewidth]{figures/calibration_acceptance.png}
  \caption{Survival probability of a candidate canary under the healthy test
  (left, for one to three healthy targets) and the weak test (right).}
  \label{fig:calibration-acceptance}
\end{figure}

\begin{table}[htbp]
  \centering
  \small
  \begin{tabular}{@{}lrrr@{\qquad}lr@{}}
    \toprule
    & \multicolumn{3}{c}{$P_H(p)$ for $|T_H| =$} & & \\
    \cmidrule(lr){2-4}
    $p$ & 1 & 2 & 3 & $w$ & $P_W(w)$ \\
    \midrule
    0.99 & 0.970 & 0.941 & 0.914 & 0.2 & 0.896 \\
    0.98 & 0.941 & 0.886 & 0.834 & 0.5 & 0.500 \\
    0.95 & 0.857 & 0.735 & 0.630 & 0.6 & 0.352 \\
         &       &       &       & 0.8 & 0.104 \\
    \bottomrule
  \end{tabular}
  \caption{Survival probabilities from \eqref{eq:acceptance}, computed by
  \texttt{figures/make\_figures.py}.}
  \label{tab:acceptance_values}
\end{table}

First, the healthy test is a conjunction over $3|T_H|$ attempts, so each
additional healthy provider lowers the survival probability geometrically. A
canary with $p = 0.98$ survives with probability $0.941$ against one provider
and $0.886$ against two. A larger healthy set yields a more conservative bank
and a smaller one.

Second, the weak test is a filter on discriminating power against one specific
weak model. A canary that the weak model answers correctly with probability
$0.6$ still survives with probability $0.352$, so the retained set contains
canaries on which the weak model is not uniformly wrong.

Third, the activation threshold can be binding. For illustration, if the
per-attempt healthy accuracy were $0.98$ on every canary, with two healthy
targets, and the weak model were correct with probability $0.5$ on every
canary, the expected number of survivors would be $300 \times 0.886 \times
0.500 \approx 133$, which is below $m = 150$. These values are assumptions
chosen to show the interaction of the thresholds and are not measurements. Real
canaries differ in difficulty, so the realized count can differ in either
direction.

The independence assumption is optimistic. All calls use temperature $0$, so
repeated attempts at one target are expected to be strongly correlated, and the
repetitions mainly check the stability of the serving path rather than provide
independent samples.

% ==============================================================================
\subsection{Runtime, persistence, and activation}
\label{sec:calibration_runtime}

The function \texttt{run\_calibration} executes the protocol and persists it in
five steps. (i) It registers the bank as \texttt{DRAFT} canaries. (ii) It
opens a run row in state \texttt{RUNNING}, with the grader version, call
parameters, thresholds, and targets stored as JSON. (iii) It collects all
observations concurrently under a semaphore of size $K$. (iv) In one
transaction it stores every raw response and grade. (v) In a second
transaction it writes each decision into the canary row (the flag, the two
accuracies, the status, and the reason) and marks the run \texttt{COMPLETED}. Any
exception marks the run \texttt{FAILED} and leaves canary statuses unchanged.

Accepted canaries move to \texttt{CALIBRATED}. Rejected canaries return to
\texttt{DRAFT}, including a canary that a previous run had accepted.
Activation is a separate step. The function \texttt{activate\_calibrated\_bank}
takes a completed run, counts the canaries that were accepted by that run, and
promotes them to \texttt{ACTIVE} only if the count reaches $m$. The run
identifier ties activation to one set of decisions, and repeating the call is
safe. The prober loads only \texttt{ACTIVE} canaries of the configured bank
version. A canary that is already \texttt{ACTIVE} cannot be updated by a
calibration run, so recalibrating requires a new bank version.
Figure~\ref{fig:lifecycles} shows the three state machines involved, including
the probe states of Section~\ref{sec:probing}.

% rendered_images:begin
% ```graphviz
% digraph Lifecycles {
%   // lifecycle states : teal | final or active : sage green | failure : coral
%   bgcolor="transparent";
%   pad="0.10";
%   splines=spline;
%   nodesep=0.25;
%   ranksep=0.38;
%   node [shape=box, style="rounded,filled", fillcolor="#FFFFFF", color="#9AA9B8", penwidth=1.8, fontname="Helvetica", fontcolor="#243B53", fontsize=11, margin="0.14,0.08", height=0.42];
%   edge [color="#A3B1C0", penwidth=1.3, arrowhead=vee, arrowsize=0.75, fontname="Helvetica", fontsize=9.5, fontcolor="#7B8794"];
%   rankdir=TB;
%   nodesep=0.35;
%   ranksep=0.5;
%   node [fillcolor="#B7DDD0", color="#6FA890", fontcolor="#1F4E39"];
% 
%   subgraph cluster_canary {
%     label="Canary status"; labelloc="t"; fontname="Helvetica-Bold"; fontsize=12; fontcolor="#3A4A5C";
%     style="rounded,filled"; fillcolor="#F7F9FB"; color="#C7D0DA"; penwidth=1.0; margin=12;
%     C_draft [label="DRAFT"];
%     C_cal   [label="CALIBRATED"];
%     C_act   [label="ACTIVE", fillcolor="#DFEDE0", color="#8FB79A", fontcolor="#2E5A3D"];
%     C_draft -> C_cal [label="Accepted"];
%     C_cal -> C_draft [label="Rejected later", style=dashed];
%     C_cal -> C_act [label="Activation"];
%   }
%   subgraph cluster_run {
%     label="Calibration run"; labelloc="t"; fontname="Helvetica-Bold"; fontsize=12; fontcolor="#3A4A5C";
%     style="rounded,filled"; fillcolor="#F7F9FB"; color="#C7D0DA"; penwidth=1.0; margin=12;
%     R_run  [label="RUNNING"];
%     R_done [label="COMPLETED", fillcolor="#DFEDE0", color="#8FB79A", fontcolor="#2E5A3D"];
%     R_fail [label="FAILED", fillcolor="#F6C6C6", color="#D98C8C", fontcolor="#6B2A2A"];
%     R_run -> R_done [label="Stored"];
%     R_run -> R_fail [label="Exception", style=dashed];
%   }
%   subgraph cluster_probe {
%     label="Probe slot"; labelloc="t"; fontname="Helvetica-Bold"; fontsize=12; fontcolor="#3A4A5C";
%     style="rounded,filled"; fillcolor="#F7F9FB"; color="#C7D0DA"; penwidth=1.0; margin=12;
%     P_sch  [label="SCHEDULED"];
%     P_run  [label="RUNNING"];
%     P_done [label="COMPLETED", fillcolor="#DFEDE0", color="#8FB79A", fontcolor="#2E5A3D"];
%     P_fail [label="FAILED_INTERNAL", fillcolor="#F6C6C6", color="#D98C8C", fontcolor="#6B2A2A"];
%     P_sch -> P_run [label="Started"];
%     P_run -> P_done [label="Linked"];
%     P_run -> P_fail [label="Exception", style=dashed];
%     P_sch -> P_fail [label="Exception", style=dashed];
%   }
% }
% ```
% label=fig:lifecycles
% caption=Lifecycles of a canary, a calibration run, and a probe slot. Teal states are in progress, green states are successful terminal or active states, and coral states are failures. Dashed edges are exceptions or later rejections.
% rendered_images:end
% render_images:begin
\begin{figure}[H]
  \includegraphics[width=\linewidth]{05_calibration.tex.figs/05_calibration.2.png}
  \caption{Lifecycles of a canary, a calibration run, and a probe slot. Teal states are in progress, green states are successful terminal or active states, and coral states are failures. Dashed edges are exceptions or later rejections.}
  \label{fig:lifecycles}
\end{figure}
% render_images:end

\subsection{Reporting}
\label{sec:calibration_report}

After a run, the summary records for each target the number of calls, the
number of successful and gradable calls, the accuracy on all gradable calls
and on the surviving canaries, the nearest-rank 95th percentile of latency, and
the cost. It also records the number of candidates, the accepted and rejected
counts, the counts per reason code, the total number of calls, and the total
cost. A markdown renderer writes these figures to
\texttt{docs/findings/calibration.md}. In the version of the code base that this
paper describes, that file is a placeholder that states that the reviewed live
run is pending and contains no measurements.
